\documentclass[11pt]{article}
\usepackage{mathblatt}


\blattfuss{Prozent}{Lösungen}
\begin{document}
\noindent{\large\bfseries Prozent \textperiodcentered{} Lösungen}\par\medskip
\begin{pfloesung}{}
\lz{1.}{$\tfrac{3}{5}$}{Bruch: $\tfrac{6}{10}$}{}
\end{pfloesung}
\begin{pfloesung}{}
\lz{2.}{$\tfrac{70}{100}$}{}{}
\end{pfloesung}
\begin{pfloesung}{}
\lz{3.}{$0{,}5$}{}{}
\end{pfloesung}
\begin{pfloesung}{}
\lz{4.}{$\tfrac{9}{10}$}{}{}
\end{pfloesung}
\begin{pfloesung}{}
\lz{5.}{$6$}{}{}
\end{pfloesung}
\begin{pfloesung}{}
\lz{6.}{$7$}{}{}
\end{pfloesung}
\begin{pfloesung}{}
\lz{7.}{$45\,€$}{Anteil: $1$ Karte: $18 : 2 = 9\,€$; Anteil: $5$ Karten: $5 \cdot 9$}{}
\end{pfloesung}
\begin{pfloesung}{}
\lz{8.}{$9\,€$}{Anteil: $1$ kg: $12 : 4 = 3\,€$; Anteil: $3$ kg: $3 \cdot 3$}{}
\end{pfloesung}
\begin{pfloesung}{}
\lz{9.}{$4{,}7$}{}{}
\end{pfloesung}
\begin{pfloesung}{}
\lz{10.}{$8{,}3$}{}{}
\end{pfloesung}
\begin{pfloesung}{}
\lz{11.}{$5\,€$}{}{}
\end{pfloesung}
\begin{pfloesung}{}
\lz{12.}{$18$ kg}{}{}
\end{pfloesung}
\begin{pfloesung}{}
\lz{13.}{64}{}{}
\end{pfloesung}
\begin{pfloesung}{}
\lz{14.}{125}{}{}
\end{pfloesung}
\begin{pfloesung}{}
\lz{15.}{jeder $5.$}{Bruch: $20\,\% = \tfrac{20}{100} = \tfrac{1}{5}$}{}
\end{pfloesung}
\begin{pfloesung}{}
\lz{16.}{$55\,\%$}{Bruch: $\tfrac{55}{100}$}{}
\end{pfloesung}
\begin{pfloesung}{}
\lz{17.}{$50\,\%$}{}{}
\end{pfloesung}
\begin{pfloesung}{}
\lz{18.}{$20\,\%$}{}{}
\end{pfloesung}
\begin{pfloesung}{{\small\color{mbgrau}P10 2022 · 1f}}
\lz{19.}{20 \%}{Bruch: $\tfrac{1}{5}$ = $\tfrac{20}{100}$}{}
\end{pfloesung}
\begin{pfloesung}{{\small\color{mbgrau}P10 2020 · 1a}}
\lz{20.}{4 von 100 Bäumen wachsen nicht an}{}{}
\end{pfloesung}
\begin{pfloesung}{}
\lz{21.}{$10$ gleiche Teile}{}{}
\end{pfloesung}
\begin{pfloesung}{}
\lz{22.}{$4$ gleiche Teile}{}{}
\end{pfloesung}
\begin{pfloesung}{}
\lz{23.}{$28$ Kinder}{}{}
\end{pfloesung}
\begin{pfloesung}{}
\lz{24.}{$600$ Lose}{}{}
\end{pfloesung}
\begin{pfloesung}{}
\lz{25.}{$12\,€$}{ein Viertel vom Ganzen: $48 : 4$}{}
\end{pfloesung}
\begin{pfloesung}{}
\lz{26.}{$16\,€$}{ein Viertel vom Ganzen: $64 : 4$}{}
\end{pfloesung}
\begin{pfloesung}{}
\lz{27.}{$48\,€$}{Rabatt $64 : 4 = 16\,€$; neuer Preis $64 - 16$}{}
\end{pfloesung}
\begin{pfloesung}{}
\lz{28.}{$0{,}90\,€$}{Prozentwert: $0{,}2 \cdot 4{,}50$}{}
\end{pfloesung}
\begin{pfloesung}{}
\lz{29.}{$2{,}80\,€$}{Prozentwert: $0{,}07 \cdot 40$}{}
\end{pfloesung}
\begin{pfloesung}{{\small\color{mbgrau}P10 2026 FOR · 1a}}
\lz{30.}{21 €}{Prozentwert: W = G · p = 70 € · 0,3 = 21 €}{}
\end{pfloesung}
\begin{pfloesung}{{\small\color{mbgrau}P10 2014 · 1a}}
\lz{31.}{6,50 €}{}{}
\end{pfloesung}
\begin{pfloesung}{}
\lz{32.}{$1{,}40\,€$, in Laden A ist die Jacke billiger}{Laden A: $0{,}8 \cdot 95 = 76\,€$; Laden B: $0{,}9 \cdot 86 = 77{,}40\,€$; Unterschied: $77{,}40 - 76$}{}
\end{pfloesung}
\begin{pfloesung}{{\small\color{mbgrau}P10 2019 · 5a}}
\lz{33.}{276 Jugendliche}{Anzahl: 23 \% von 1200 = 276}{}
\end{pfloesung}
\begin{pfloesung}{{\small\color{mbgrau}P10 2017 · 1b}}
\lz{34.}{24 €}{}{}
\end{pfloesung}
\begin{pfloesung}{{\small\color{mbgrau}P10 2021 · 1c}}
\lz{35.}{110 €}{}{}
\end{pfloesung}
\begin{pfloesung}{{\small\color{mbgrau}P10 2015 · 2a}}
\lz{36.}{Term 1 falsch; Term 2 richtig; Term 3 richtig}{voller Preis 3 · 12 € + 7,50 € = 43,50 €; Rabatt 8,70 €}{}
\end{pfloesung}
\begin{pfloesung}{}
\lz{37.}{$5$ gleiche Teile}{}{}
\end{pfloesung}
\begin{pfloesung}{}
\lz{38.}{$20$ gleiche Teile}{Prozentsatz: je $5$ mm}{}
\end{pfloesung}
\begin{pfloesung}{}
\lz{39.}{$150$ Gäste}{das Ganze steht hinten}{}
\end{pfloesung}
\begin{pfloesung}{}
\lz{40.}{der alte Preis, $60\,€$}{}{}
\end{pfloesung}
\begin{pfloesung}{}
\lz{41.}{$20\,\%$}{}{}
\end{pfloesung}
\begin{pfloesung}{}
\lz{42.}{$75\,\%$}{}{}
\end{pfloesung}
\begin{pfloesung}{}
\lz{43.}{$30\,\%$}{Ganzes $20$; Prozentsatz: $6 : 20 = 0{,}3$}{}
\end{pfloesung}
\begin{pfloesung}{}
\lz{44.}{$25\,\%$}{Rabatt $20\,€$; Prozentsatz: $20 : 80 = 0{,}25$}{}
\end{pfloesung}
\begin{pfloesung}{}
\lz{45.}{$17\,\%$}{Hundertstel: $\tfrac{17}{100}$}{}
\end{pfloesung}
\begin{pfloesung}{}
\lz{46.}{$68\,\%$}{Hundertstel: $\tfrac{68}{100}$}{}
\end{pfloesung}
\begin{pfloesung}{}
\lz{47.}{am Dienstag}{Montag $36 : 240 = 0{,}15 = 15\,\%$; Dienstag $27 : 150 = 0{,}18 = 18\,\%$}{}
\end{pfloesung}
\begin{pfloesung}{}
\lz{48.}{$\tfrac{17}{24} \approx 70{,}8\,\%$}{Prozentsatz: $17 : 24 = 0{,}7083\ldots$}{}
\end{pfloesung}
\begin{pfloesung}{{\small\color{mbgrau}P10 2018 · 7a}}
\lz{49.}{12,5 \%}{}{}
\end{pfloesung}
\begin{pfloesung}{{\small\color{mbgrau}P10 2015 · 7c}}
\lz{50.}{$\tfrac{12}{83}$ $\approx$ 14,5 \%; Sektor $\approx$ 52°; $\approx$ 52°}{Anzahl: 83 $-$ 71 = 12; Anteil: 12 : 83 = $\tfrac{12}{83}$ $\approx$ 0,145; Winkel: $\tfrac{12}{83}$ · 360° = $\tfrac{4320}{83}$° $\approx$ 52°}{}
\end{pfloesung}
\begin{pfloesung}{}
\lz{51.}{$10$ Abschnitte bis $100\,\%$}{}{}
\end{pfloesung}
\begin{pfloesung}{}
\lz{52.}{$2$ Abschnitte bis $100\,\%$}{}{}
\end{pfloesung}
\begin{pfloesung}{}
\lz{53.}{Prozentsatz}{gegeben: Grundwert $240$, Prozentwert $60$}{}
\end{pfloesung}
\begin{pfloesung}{}
\lz{54.}{Grundwert}{gegeben: Prozentsatz $30\,\%$, Prozentwert $9$ Kinder}{}
\end{pfloesung}
\begin{pfloesung}{}
\lz{55.}{$240$ kg}{wie oft bis hundert Prozent: $100 : 10 = 10$; mal nehmen: $10 \cdot 24 = 240$ kg}{}
\end{pfloesung}
\begin{pfloesung}{}
\lz{56.}{$48$ kg}{wie oft bis hundert Prozent: $100 : 50 = 2$; mal nehmen: $2 \cdot 24 = 48$ kg}{}
\end{pfloesung}
\begin{pfloesung}{}
\lz{57.}{$60$ l}{Grundwert: $5 \cdot 12$}{}
\end{pfloesung}
\begin{pfloesung}{}
\lz{58.}{$225$ kg}{Grundwert: $81 : 36 = 2{,}25$; Grundwert: $2{,}25 \cdot 100$}{}
\end{pfloesung}
\begin{pfloesung}{}
\lz{59.}{Ja}{ein Prozent: $7 : 28 = 0{,}25$ h; hundert Prozent: $100 \cdot 0{,}25 = 25$ Stunden}{}
\end{pfloesung}
\begin{pfloesung}{{\small\color{mbgrau}P10 2025 · 1a}}
\lz{60.}{30 €}{Grundwert: G = W : p = 6 € : 0,2 = 30 €}{}
\end{pfloesung}
\begin{pfloesung}{{\small\color{mbgrau}P10 2023 · 1b}}
\lz{61.}{800 €}{Grundwert: G = W : p = 200 € : 0,25 = 800 €}{}
\end{pfloesung}
\begin{pfloesung}{}
\lz{62.}{auf $80\,\%$ gesenkt}{}{}
\end{pfloesung}
\begin{pfloesung}{}
\lz{63.}{$120$ Mitglieder}{steht hinten}{}
\end{pfloesung}
\begin{pfloesung}{}
\lz{64.}{$600\,€$}{}{}
\end{pfloesung}
\begin{pfloesung}{}
\lz{65.}{um $40\,\%$ gesenkt}{}{}
\end{pfloesung}
\begin{pfloesung}{}
\lz{66.}{$77\,€$}{Prozentwert: $10\,\% = 7\,€$; dazu: $70 + 7 = 77\,€$}{}
\end{pfloesung}
\begin{pfloesung}{}
\lz{67.}{$60\,€$}{Anteil: $66 : 1{,}1$}{}
\end{pfloesung}
\begin{pfloesung}{}
\lz{68.}{$105\,€$}{Prozentwert: $50\,\% = 35\,€$; dazu: $70 + 35 = 105\,€$}{}
\end{pfloesung}
\begin{pfloesung}{}
\lz{69.}{$25\,\%$}{Anteil: $10 : 40 = 0{,}25$}{}
\end{pfloesung}
\begin{pfloesung}{}
\lz{70.}{$50\,€$}{Anteil: $53{,}50 : 1{,}07$}{}
\end{pfloesung}
\begin{pfloesung}{}
\lz{71.}{$15\,\%$}{Veränderung gesucht: $(4{,}83 - 4{,}20) : 4{,}20 = 0{,}15$}{}
\end{pfloesung}
\begin{pfloesung}{}
\lz{72.}{Ja}{erster Rabatt: $600 \cdot 0{,}9 = 540\,€$; zweiter Rabatt: $540 \cdot 0{,}9 = 486\,€$; vergleichen: $486\,€ < 490\,€$; Rest: $490 - 486 = 4\,€$}{}
\end{pfloesung}
\begin{pfloesung}{{\small\color{mbgrau}P10 2020 · 2d}}
\lz{73.}{280 $-$ 130 = 150; 210 + 210 : 3 = 280}{ein Drittel von 210 = 70}{}
\end{pfloesung}
\begin{pfloesung}{{\small\color{mbgrau}P10 2020 · 4c}}
\lz{74.}{nein, Studie 1 54 · 1,03$^{14}$ $\approx$ 81,7 gegen Studie 2 = 75,6 (Tausend Liter)}{Verbrauch: 54 · 1,03$^{14}$ $\approx$ 81,7; Verbrauch: 54 · 1,4 = 75,6}{}
\end{pfloesung}
\begin{pfloesung}{}
\lz{75.}{$50\,\%$ mehr}{Rechnung: $1$ Prozentpunkt; Anteil: $1 : 2 = 0{,}5$}{}
\end{pfloesung}
\begin{pfloesung}{{\small\color{mbgrau}P10 2026 FOR · 3c}}
\lz{76.}{$\tfrac{5}{179}$ $\approx$ 2,8 \%}{Differenz 0,05 €; Anteil: 0,05 : 1,79 = $\tfrac{5}{179}$ $\approx$ 0,0279}{}
\end{pfloesung}
\begin{pfloesung}{{\small\color{mbgrau}P10 2024 · 1e}}
\lz{77.}{4,20 €}{Unterschied: 20 \% von 3,50 € = 0,70 €}{}
\end{pfloesung}
\begin{pfloesung}{{\small\color{mbgrau}P10 2016 · 2c}}
\lz{78.}{25 \%}{Anzahl: 400 000 $-$ 320 000 = 80 000}{}
\end{pfloesung}
\begin{pfloesung}{{\small\color{mbgrau}P10 2022 · 4b}}
\lz{79.}{$\tfrac{7}{47}$ $\approx$ 14,9 \%}{Differenz 70; Anteil: 70 : 470 = $\tfrac{7}{47}$ $\approx$ 0,149}{}
\end{pfloesung}
\begin{pfloesung}{{\small\color{mbgrau}P10 2015 · 2b}}
\lz{80.}{34,80 €}{Prozent: 43,50 € · 0,8 = 34,80 €}{}
\end{pfloesung}
\begin{pfloesung}{{\small\color{mbgrau}P10 2021 · 5a}}
\lz{81.}{2,7 Mio.; 8,7 Mio.; $\tfrac{11}{79}$ $\approx$ 13,9 \%}{Unterschied: 9,7 $-$ 7,0 = 2,7; Summe 78,3 Mio.; Summe: 78,3 : 9 = 8,7; Anzahl: 9,0 $-$ 7,9 = 1,1 Mio.}{}
\end{pfloesung}
\begin{pfloesung}{}
\lz{82.}{$2{,}40\,€$}{}{}
\end{pfloesung}
\begin{pfloesung}{}
\lz{83.}{$425\,€$}{der alte ist hier der größere}{}
\end{pfloesung}
\begin{pfloesung}{}
\lz{84.}{auf $115\,\%$ erhöht}{}{}
\end{pfloesung}
\begin{pfloesung}{}
\lz{85.}{um $25\,\%$ erhöht}{}{}
\end{pfloesung}
\begin{pfloesung}{}
\lz{86.}{$52{,}50\,€$}{Prozentwert: $25\,\% = 17{,}50\,€$; weg: $70 - 17{,}50 = 52{,}50\,€$}{}
\end{pfloesung}
\begin{pfloesung}{}
\lz{87.}{$56\,€$}{Prozentwert: $20\,\% = 14\,€$; weg: $70 - 14 = 56\,€$}{}
\end{pfloesung}
\begin{pfloesung}{}
\lz{88.}{$7\,\%$}{Anteil: $14 : 200 = 0{,}07$}{}
\end{pfloesung}
\begin{pfloesung}{}
\lz{89.}{$70\,\%$}{}{}
\end{pfloesung}
\begin{pfloesung}{}
\lz{90.}{$30\,\%$}{}{}
\end{pfloesung}
\begin{pfloesung}{}
\lz{91.}{in der 7a ist der Anteil größer}{Bruch: 7a: $\tfrac{18}{25} = 72\,\%$; Bruch: 7b: $\tfrac{21}{30} = \tfrac{7}{10} = 70\,\%$}{}
\end{pfloesung}
\begin{pfloesung}{{\small\color{mbgrau}P10 2019 · 5b}}
\lz{92.}{Aussage 1 falsch; etwa vier Fünftel (81 \%) der Mädchen nutzen einen Fernseher}{drei Viertel = 75 \%}{}
\end{pfloesung}
\begin{pfloesung}{{\small\color{mbgrau}P10 2025 · 4b}}
\lz{93.}{100 m und 6 m; Ja, Steigung 16 : 170 = $\tfrac{8}{85}$ $\approx$ 9,4 \% > 6 \%}{Verhältnis: 16 : 170 = $\tfrac{8}{85}$ $\approx$ 0,0941}{}
\end{pfloesung}
\begin{pfloesung}{{\small\color{mbgrau}P10 2017 · 2b}}
\lz{94.}{34,4 \% + 11,2 \% + 5,0 \% = 50,6 \% > 50 \%}{}{}
\end{pfloesung}
\begin{pfloesung}{}
\lz{95.}{$\approx 21{,}5\,\%$}{Anteil: $40 : 10 = 4$ Untersetzer je Reihe, $4 \cdot 4 = 16$ Untersetzer; Untersetzer $16 \cdot \pi \cdot 5^2 \approx 1\,256{,}6\,\text{cm}^2$; Platte $40^2 = 1\,600\,\text{cm}^2$; Abfall $1\,600 - 1\,256{,}6 = 343{,}4\,\text{cm}^2$; Anteil: $343{,}4 : 1\,600$}{}
\end{pfloesung}
\begin{pfloesung}{{\small\color{mbgrau}P10 2023 · 5c}}
\lz{96.}{Ja, 12 Deckel, Abfall 800 $-$ 192$\pi$ $\approx$ 196,8 cm² $\approx$ 24,6 \% $\le$ 25 \%}{Anzahl: 100 : 8 = 12,5 $\Rightarrow$ 12 Deckel; Deckelfläche 12 · $\pi$ · 16 = 192$\pi$ $\approx$ 603,2 cm²; Blech 800 cm²}{}
\end{pfloesung}
\begin{pfloesung}{}
\lz{97.}{der Zinssatz}{Prozentsatz}{}
\end{pfloesung}
\begin{pfloesung}{}
\lz{98.}{$K \cdot p\,\%$}{Zinssatz: $Z$; Zinsen = Prozentwert, Kapital = Grundwert}{}
\end{pfloesung}
\begin{pfloesung}{}
\lz{99.}{3\,\%}{Zinssatz: $27 : 900 = 0{,}03$}{}
\end{pfloesung}
\begin{pfloesung}{}
\lz{100.}{1,5\,\%}{Zinssatz: $60 : 4\,000 = 0{,}015$}{}
\end{pfloesung}
\begin{pfloesung}{}
\lz{101.}{700\,€}{Zinssatz: 1\,\% = 7\,€}{}
\end{pfloesung}
\begin{pfloesung}{}
\lz{102.}{stimmt}{Zinssatz: $300 \cdot 0{,}03 = 9{,}00$\,€; oder $9 : 300 = 0{,}03 = 3\,\%$}{}
\end{pfloesung}
\begin{pfloesung}{}
\lz{103.}{6\,€}{Zinssatz: $300 \cdot 0{,}02$}{}
\end{pfloesung}
\begin{pfloesung}{}
\lz{104.}{21\,€}{Zinssatz: $700 \cdot 0{,}03$}{}
\end{pfloesung}
\begin{pfloesung}{{\small\color{mbgrau}P10 2015 · 1e}}
\lz{105.}{8 €}{}{}
\end{pfloesung}
\begin{pfloesung}{{\small\color{mbgrau}P10 2014 · 1e}}
\lz{106.}{2,3 \%}{}{}
\end{pfloesung}
\begin{pfloesung}{{\small\color{mbgrau}P10 2014 · 3a}}
\lz{107.}{200 € · 0,02 = 4,00 €}{}{}
\end{pfloesung}
\begin{pfloesung}{}
\lz{108.}{vom neuen Guthaben 612\,€}{}{}
\end{pfloesung}
\begin{pfloesung}{}
\lz{109.}{vom neuen Guthaben 2\,060\,€}{}{}
\end{pfloesung}
\begin{pfloesung}{}
\lz{110.}{5\,050\,€}{Zinsen 50\,€}{}
\end{pfloesung}
\begin{pfloesung}{}
\lz{111.}{2\,310\,€}{Zinsen 110\,€}{}
\end{pfloesung}
\begin{pfloesung}{}
\lz{112.}{Guthaben = 477,00\,€, Zinsen = 9,54\,€}{}{}
\end{pfloesung}
\begin{pfloesung}{}
\lz{113.}{1\,003,97\,€}{Guthaben: 732,00\,€; Zinseszins: 866,64\,€; Zinseszins: 1\,003,97\,€}{}
\end{pfloesung}
\begin{pfloesung}{}
\lz{114.}{$\approx$ 3\,939,28\,€}{Zinseszins: $3\,500 \cdot 1{,}03^{4}$}{}
\end{pfloesung}
\begin{pfloesung}{{\small\color{mbgrau}P10 2014 · 3b}}
\lz{115.}{824,32 €; 16,49 €}{}{}
\end{pfloesung}
\begin{pfloesung}{{\small\color{mbgrau}P10 2014 · 3c}}
\lz{116.}{1000 € · 1,03$^{5}$ $\approx$ 1159,27 €}{Wachstumsfaktor 1,03}{}
\end{pfloesung}
\begin{pfloesung}{}
\lz{117.}{Ja, knapp: nach 4 Jahren $\approx$ 2\,251,02\,€}{}{}
\end{pfloesung}
\begin{pfloesung}{\textbf{Prüfstein} \textperiodcentered{} P10 2023 · 6}
\lz{a)}{$\tfrac{25}{156}$ $\approx$ 16,0 \%}{Anteil: 1,25 : 7,8 = $\tfrac{25}{156}$ $\approx$ 0,160}{}
\lz{b)}{Aussage 2 falsch; Ein Viertel der Weltbevölkerung waren 2020 Kinder (1,95 : 7,8 = 0,25)}{mittleres Alter 3,9 : 7,8 = 50 \%; Kinder 25 \%}{}
\end{pfloesung}
\end{document}